\documentclass[10pt,a4paper]{amsart}
\usepackage{amsmath,amssymb,amsthm,graphicx}
\usepackage{verbatim}
\newenvironment{reference}{\comment}{\endcomment}
%\newenvironment{reference}{}{}
\newenvironment{slogan}{\comment}{\endcomment}
\newenvironment{history}{\comment}{\endcomment}

% For commutative diagrams we use Xy-pic
\usepackage[all]{xy}

% We use 2cell for 2-commutative diagrams.
\xyoption{2cell}
% Theorem environments.
%
\theoremstyle{plain}
\newtheorem{theorem}[subsection]{Theorem}
\newtheorem{proposition}[subsection]{Proposition}
\newtheorem{lemma}[subsection]{Lemma}

\theoremstyle{definition}
\newtheorem{definition}[subsection]{Definition}
\newtheorem{example}[subsection]{Example}
\newtheorem{exercise}[subsection]{Exercise}
\newtheorem{situation}[subsection]{Situation}

\theoremstyle{remark}
\newtheorem{remark}[subsection]{Remark}
\newtheorem{remarks}[subsection]{Remarks}

\numberwithin{equation}{subsection}

% Macros
%
\def\lim{\mathop{\mathrm{lim}}\nolimits}
\def\colim{\mathop{\mathrm{colim}}\nolimits}
\def\Spec{\mathop{\mathrm{Spec}}}
\def\Hom{\mathop{\mathrm{Hom}}\nolimits}
\def\Ext{\mathop{\mathrm{Ext}}\nolimits}
\def\SheafHom{\mathop{\mathcal{H}\!\mathit{om}}\nolimits}
\def\SheafExt{\mathop{\mathcal{E}\!\mathit{xt}}\nolimits}
\def\Sch{\mathit{Sch}}
\def\Mor{\mathop{\mathrm{Mor}}\nolimits}
\def\Ob{\mathop{\mathrm{Ob}}\nolimits}
\def\Sh{\mathop{\mathit{Sh}}\nolimits}
\def\NL{\mathop{N\!L}\nolimits}
\def\CH{\mathop{\mathrm{CH}}\nolimits}
\def\proetale{{pro\text{-}\acute{e}tale}}
\def\etale{{\acute{e}tale}}
\def\QCoh{\mathit{QCoh}}
\def\Ker{\mathop{\mathrm{Ker}}}
\def\Im{\mathop{\mathrm{Im}}}
\def\Coker{\mathop{\mathrm{Coker}}}
\def\Coim{\mathop{\mathrm{Coim}}}

% Boxtimes
%
\DeclareMathSymbol{\boxtimes}{\mathbin}{AMSa}{"02}

%
% Macros for moduli stacks/spaces
%
\def\QCohstack{\mathcal{QC}\!\mathit{oh}}
\def\Cohstack{\mathcal{C}\!\mathit{oh}}
\def\Spacesstack{\mathcal{S}\!\mathit{paces}}
\def\Quotfunctor{\mathrm{Quot}}
\def\Hilbfunctor{\mathrm{Hilb}}
\def\Curvesstack{\mathcal{C}\!\mathit{urves}}
\def\Polarizedstack{\mathcal{P}\!\mathit{olarized}}
\def\Complexesstack{\mathcal{C}\!\mathit{omplexes}}
% \Pic is the operator that assigns to X its picard group, usage \Pic(X)
% \Picardstack_{X/B} denotes the Picard stack of X over B
% \Picardfunctor_{X/B} denotes the Picard functor of X over B
\def\Pic{\mathop{\mathrm{Pic}}\nolimits}
\def\Picardstack{\mathcal{P}\!\mathit{ic}}
\def\Picardfunctor{\mathrm{Pic}}
\def\Deformationcategory{\mathcal{D}\!\mathit{ef}}
\usepackage{cleveref}
\crefname{originalreference}{original reference}{original references}
\Crefname{originalreference}{Original reference}{Original references}
\makeatletter
\newlabel{definition-multicross}{{0C1W}{source}{}{sourcepointer-1}{}}
\newlabel{definition-multicross@cref}{{[originalreference][0][]0C1W}{[1][0][]source}{}{}{}}
\makeatother
\crefname{originaldefinition}{original definition}{original definitions}
\Crefname{originaldefinition}{Original definition}{Original definitions}
\makeatletter
\newlabel{duality-definition-gorenstein}{{0AWW}{source}{}{sourcepointer-2}{}}
\newlabel{duality-definition-gorenstein@cref}{{[originaldefinition][0][]0AWW}{[1][0][]source}{}{}{}}
\makeatother
\crefname{originallemma}{original lemma}{original lemmas}
\Crefname{originallemma}{Original lemma}{Original lemmas}
\makeatletter
\newlabel{duality-lemma-gorenstein-lci}{{0BVA}{source}{}{sourcepointer-3}{}}
\newlabel{duality-lemma-gorenstein-lci@cref}{{[originallemma][0][]0BVA}{[1][0][]source}{}{}{}}
\makeatother
\crefname{originallemma}{original lemma}{original lemmas}
\Crefname{originallemma}{Original lemma}{Original lemmas}
\makeatletter
\newlabel{dualizing-lemma-find-function}{{0A7L}{source}{}{sourcepointer-4}{}}
\newlabel{dualizing-lemma-find-function@cref}{{[originallemma][0][]0A7L}{[1][0][]source}{}{}{}}
\makeatother
\crefname{originallemma}{original lemma}{original lemmas}
\Crefname{originallemma}{Original lemma}{Original lemmas}
\makeatletter
\newlabel{dualizing-lemma-flat-under-gorenstein}{{0BJL}{source}{}{sourcepointer-5}{}}
\newlabel{dualizing-lemma-flat-under-gorenstein@cref}{{[originallemma][0][]0BJL}{[1][0][]source}{}{}{}}
\makeatother
\crefname{originallemma}{original lemma}{original lemmas}
\Crefname{originallemma}{Original lemma}{Original lemmas}
\makeatletter
\newlabel{dualizing-lemma-regular-gorenstein}{{0AWX}{source}{}{sourcepointer-6}{}}
\newlabel{dualizing-lemma-regular-gorenstein@cref}{{[originallemma][0][]0AWX}{[1][0][]source}{}{}{}}
\makeatother
\crefname{originalequation}{original equation}{original equations}
\Crefname{originalequation}{Original equation}{Original equations}
\makeatletter
\newlabel{equation-multicross}{{0C1U}{source}{}{sourcepointer-7}{}}
\newlabel{equation-multicross@cref}{{[originalequation][0][]0C1U}{[1][0][]source}{}{}{}}
\makeatother
\crefname{originalsection}{original section}{original sections}
\Crefname{originalsection}{Original section}{Original sections}
\makeatletter
\newlabel{more-algebra-section-permanence-completion}{{07NU}{source}{}{sourcepointer-8}{}}
\newlabel{more-algebra-section-permanence-completion@cref}{{[originalsection][0][]07NU}{[1][0][]source}{}{}{}}
\makeatother
\title{A Gorenstein multicross singularity over an algebraically closed field is a node}
\author{The Stacks Project Authors}
\date{stacks-project:curves.tex:local-snapshot-20261012; GNU FDL 1.2 or later}
\begin{document}
\maketitle
\noindent\textit{Editorial note.} All delivered mathematical text is the author's own, copied from the pinned Stacks Project chapter file curves.tex; no statement or proof step was rewritten, reordered or removed.\par\medskip
\noindent\textit{Editorial note.} The counted proof is the author's complete argument: the completion reduction, the n = 2 complete-intersection case, the full Ext\textasciicircum{}1 computation with the coordinate functionals e\_i, the independence check, and the final dimension contradiction are all retained verbatim.\par\medskip
\noindent\textit{Editorial note.} The label lemma-multicross-gorenstein-is-nodal names the counted statement, and the proof environment immediately following it in the source is delivered verbatim as its proof.\par\medskip
\noindent\textit{Editorial note.} The counted statement is carried by the author's own lemma environment in the chapter source.\par\medskip
\noindent\textit{Allowed standard prerequisites (editorial).}\begin{itemize}
\item The multicross normal form. Same chapter, Definition definition-multicross and Equation equation-multicross. Provides the completed local normal form A with n branches whose Gorenstein property is analysed.
\item Gorenstein permanence under completion. More on Algebra, Section more-algebra-section-permanence-completion; Dualizing Complexes, Lemma dualizing-lemma-flat-under-gorenstein. The completed local map is flat and unramified with trivial residue extension, transferring the Gorenstein property from O\_\{X,x\} to its completion.
\item Duality inputs. Duality for Schemes, Definition duality-definition-gorenstein, Lemma duality-lemma-gorenstein-lci, Lemma duality-lemma-find-function; Dualizing Complexes, Lemma dualizing-lemma-regular-gorenstein. Handles the n = 2 case (complete intersection) and, for n > 2, the constraint R Hom(k, A) = k[n] on a Gorenstein ring.
\end{itemize}\medskip
\ifdefined\hypertarget\hypertarget{sourcepointer-1}{}\fi
\noindent\textit{Source reference.} Original reference 0C1W: sources/stacks-curves/pinned-original-source.tex, source lines 3463--3463. Prerequisite pointer invoked at proof line(s) 3538; it names a standard fact of the same chapter that the counted argument uses, and every step of the counted proof built on it is retained verbatim in the delivered body.\par\medskip
\ifdefined\hypertarget\hypertarget{sourcepointer-2}{}\fi
\noindent\textit{Source reference.} Original definition 0AWW: sources/stacks-chow/tags.txt, source lines 13734--13734. Prerequisite pointer invoked at proof line(s) 3565; it names a standard fact of a sibling chapter that the counted argument uses, and every step of the counted proof built on it is retained verbatim in the delivered body.\par\medskip
\ifdefined\hypertarget\hypertarget{sourcepointer-3}{}\fi
\noindent\textit{Source reference.} Original lemma 0BVA: sources/stacks-chow/tags.txt, source lines 14903--14903. Prerequisite pointer invoked at proof line(s) 3557; it names a standard fact of a sibling chapter that the counted argument uses, and every step of the counted proof built on it is retained verbatim in the delivered body.\par\medskip
\ifdefined\hypertarget\hypertarget{sourcepointer-4}{}\fi
\noindent\textit{Source reference.} Original lemma 0A7L: sources/stacks-chow/tags.txt, source lines 12884--12884. Prerequisite pointer invoked at proof line(s) 3567; it names a standard fact of a sibling chapter that the counted argument uses, and every step of the counted proof built on it is retained verbatim in the delivered body.\par\medskip
\ifdefined\hypertarget\hypertarget{sourcepointer-5}{}\fi
\noindent\textit{Source reference.} Original lemma 0BJL: sources/stacks-chow/tags.txt, source lines 14529--14529. Prerequisite pointer invoked at proof line(s) 3549; it names a standard fact of a sibling chapter that the counted argument uses, and every step of the counted proof built on it is retained verbatim in the delivered body.\par\medskip
\ifdefined\hypertarget\hypertarget{sourcepointer-6}{}\fi
\noindent\textit{Source reference.} Original lemma 0AWX: sources/stacks-chow/tags.txt, source lines 13735--13735. Prerequisite pointer invoked at proof line(s) 3560; it names a standard fact of a sibling chapter that the counted argument uses, and every step of the counted proof built on it is retained verbatim in the delivered body.\par\medskip
\ifdefined\hypertarget\hypertarget{sourcepointer-7}{}\fi
\noindent\textit{Source reference.} Original equation 0C1U: sources/stacks-curves/pinned-original-source.tex, source lines 3410--3410. Prerequisite pointer invoked at proof line(s) 3551; it names a standard fact of the same chapter that the counted argument uses, and every step of the counted proof built on it is retained verbatim in the delivered body.\par\medskip
\ifdefined\hypertarget\hypertarget{sourcepointer-8}{}\fi
\noindent\textit{Source reference.} Original section 07NU: sources/stacks-chow/tags.txt, source lines 9774--9774. Prerequisite pointer invoked at proof line(s) 3546; it names a standard fact of a sibling chapter that the counted argument uses, and every step of the counted proof built on it is retained verbatim in the delivered body.\par\medskip
\begin{lemma}
\label{lemma-multicross-gorenstein-is-nodal}
Let $k$ be an algebraically closed field. Let $X$ be a reduced algebraic
$1$-dimensional $k$-scheme. Let $x \in X$ be a multicross singularity
(Definition \ref{definition-multicross}).
If $X$ is Gorenstein, then $x$ is a node.
\end{lemma}

\begin{proof}
The map $\mathcal{O}_{X, x} \to \mathcal{O}_{X, x}^\wedge$
is flat and unramified in the sense that
$\kappa(x) = \mathcal{O}_{X, x}^\wedge/\mathfrak m_x \mathcal{O}_{X, x}^\wedge$.
(See More on Algebra, Section \ref{more-algebra-section-permanence-completion}.)
Thus $X$ is Gorenstein implies $\mathcal{O}_{X, x}$ is Gorenstein, implies
$\mathcal{O}_{X, x}^\wedge$ is Gorenstein by
Dualizing Complexes, Lemma \ref{dualizing-lemma-flat-under-gorenstein}.
Thus it suffices to show that the ring $A$ in
(\ref{equation-multicross}) with $n \geq 2$
is Gorenstein if and only if $n = 2$.

\medskip\noindent
If $n = 2$, then $A = k[[x, y]]/(xy)$ is a complete intersection
and hence Gorenstein. For example this follows from
Duality for Schemes, Lemma \ref{duality-lemma-gorenstein-lci}
applied to $k[[x, y]] \to A$ and the fact that the regular
local ring $k[[x, y]]$ is Gorenstein by
Dualizing Complexes, Lemma \ref{dualizing-lemma-regular-gorenstein}.

\medskip\noindent
Assume $n > 2$. If $A$ where Gorenstein, then $A$ would be a
dualizing complex over $A$
(Duality for Schemes, Definition \ref{duality-definition-gorenstein}).
Then $R\Hom(k, A)$ would be equal to $k[n]$ for some $n \in \mathbf{Z}$, see
Dualizing Complexes, Lemma \ref{dualizing-lemma-find-function}.
It would follow that $\Ext^1_A(k, A) \cong k$
or $\Ext^1_A(k, A) = 0$ (depending on the value of $n$;
in fact $n$ has to be $-1$ but it doesn't matter to us here).
Using the exact sequence
$$
0 \to \mathfrak m_A \to A \to k \to 0
$$
we find that
$$
\Ext^1_A(k, A) = \Hom_A(\mathfrak m_A, A)/A
$$
where $A \to \Hom_A(\mathfrak m_A, A)$ is given by
$a \mapsto (a' \mapsto aa')$. Let $e_i \in \Hom_A(\mathfrak m_A, A)$
be the element that sends $(f_1, \ldots, f_n) \in \mathfrak m_A$ to
$(0, \ldots, 0, f_i, 0, \ldots, 0)$. The reader verifies easily
that $e_1, \ldots, e_{n - 1}$ are $k$-linearly independent in
$\Hom_A(\mathfrak m_A, A)/A$. Thus
$\dim_k \Ext^1_A(k, A) \geq n - 1 \geq 2$ which
finishes the proof.
(Observe that $e_1 + \ldots + e_n$ is the image of $1$ under the map
$A \to \Hom_A(\mathfrak m_A, A)$.)
\end{proof}

\end{document}
